\documentclass[10pt,a4paper]{amsart}
\usepackage[margin=19mm]{geometry}
\usepackage{amsmath,amssymb,mathtools}
\usepackage[scr=rsfs,cal=euler]{mathalfa}
\usepackage{xfrac}
\usepackage[unicode,hidelinks,bookmarks=false]{hyperref}
\newcommand{\ie}{i.e.\ }
\newcommand{\cf}{cf.\ }
\newcommand{\resp}{respectively\ }
\newcommand{\Subsection}{\subsection}
\providecommand{\nobreakdash}{\nobreak\hbox{-}}
\setlength{\emergencystretch}{2em}
\multlinegap0pt
\usepackage{bm}
\usepackage{bbm}
\usepackage{dsfont}
\usepackage{xspace}
\usepackage{cancel}
\usepackage{mathtools}
\usepackage{amsthm}
\makeatletter
\@ifundefined{lem}{\newtheorem{lem}{Lemma}}{}
\@ifundefined{thm}{\newtheorem{thm}{Theorem}}{}
\makeatother
\newcommand{\A}{\mathbb{A}}
\newcommand{\ZZ}{\mathbb{Z}}
\title[Formal Abel relations for curves in characteristic $p$]{Formal Abel relations for curves in characteristic $p$}
\author{John B. Little}
\date{Source: arXiv:2607.09471v1 (CC BY 4.0)}
\begin{document}
\maketitle

\noindent\emph{Source and licence.} Formal Abel relations for curves in characteristic $p$. Version arXiv:2607.09471v1 (CC BY 4.0), distributed under the Creative Commons Attribution 4.0 International licence, as recorded in the arXiv licence metadata for this exact version and shown on the abstract page https://arxiv.org/abs/2607.09471v1. Full rights notice: licenses/arxiv-2607.09471-CC-BY-4.0.txt.\par\smallskip

\noindent\textit{Editorial scope.} Counted unit: the Thm whose source label is \texttt{None} in the article (source0.tex lines 435--442, source label \texttt{None}), delivered together with the article's own complete original proof of it (source0.tex lines 444--662, one \texttt{proof} environment of 219 lines). No mathematics is written, repaired or re-derived: statement and proof are the article's own text line for line. Required context retained uncounted: FGNF1 (lines 341--348). Every definition, display and label of those lines is retained unchanged. Omitted from this package and not redistributed: 1--340, 349--434, 443--443, 663--1438. Nothing inside a retained excerpt is omitted or altered: every retained body is delivered exactly as the article's lines are written, so no omission receipt of any kind accompanies this package. Citations: the retained excerpts carry no citation command at all, so no bibliography entry of the article is transcribed and no reference is redistributed. Reference adaptation: none was needed and none was made; every reference inside the delivered unit resolves inside it, so no unresolved reference or citation is produced. Printed slips kept literal: no printed word, symbol, hypothesis or number of the counted statement or of its proof was altered. Layout and class: the article's own class is \texttt{amsart} with packages including amsmath, amsthm, amssymb, hyperref, color, setspace, graphicx, caption, subcaption, mathtools, enumerate, stackengine, dsfont, xypic; the wrapper is the lane's pinned \texttt{amsart} editorial wrapper at 10pt on A4, which restates the theorem-like environments the retained text needs and adds the article's own macro definitions that the retained text needs (\texttt{\textbackslash A}, \texttt{\textbackslash ZZ}), with extra wrapper packages (bm, bbm, dsfont, xspace, cancel, mathtools). The delivered numbering is the unit's own. Assets: no retained span contains a figure environment, a graphics-inclusion command, a PDF-page inclusion, a table or a TikZ picture; no image file is copied into this package, the package declares no asset dependency and no image file is redistributed with it. Licence: version arXiv:2607.09471v1 of this article is distributed under the Creative Commons Attribution 4.0 International licence, as recorded in the arXiv licence metadata for this exact version and shown on the abstract page https://arxiv.org/abs/2607.09471v1. A full rights notice is delivered at licenses/arxiv-2607.09471-CC-BY-4.0.txt. Characters: every retained excerpt is pure ASCII, so the delivered engine, XeLaTeX with its default Latin Modern text font, reports no missing character.
\begin{lem}
\label{FGNF1}
 Every formal group law over an algebraically closed field $k$ of characteristic $p$ is strictly isomorphic to one of the form
$$H(X,Y) \equiv X + Y \bmod \deg p.$$
If $k$ has characteristic zero, then every $m$-dimensional formal group law is strictly isomorphic to the additive $m$-dimensional 
group law
$$H(X,Y) = X + Y.$$
\end{lem}

\begin{thm}
Let $p > d + 2$, let $k$ be an algebraically closed field of characteristic $p$, 
and consider the projection $\mathcal{A}_{d+2}^{\rm max} \to \mathcal{X}_{d+2}^{\rm max}$.  If $(x_i, t_i, G) \in \mathcal{A}_{d+2}(k)$
for a particular set of mappings $t_i$ to be described below (see \eqref{Ai1s13}, \eqref{Ai2s2} below), then
the corresponding $f_i(x)$ are the local expansions of an algebraic curve of degree $d$ defined over $k$.  In other words, 
$$\mathcal{X}_{d+2}^{\rm max} =
 \{(x_i) : f_i \bmod \deg\, (d+1) \text{ from an algebraic curve of degree $d$ }\}.$$
\end{thm}

\begin{proof}
For simplicity of presentation, the details will be given for the case $d = 4$, so we assume $p \ge 7$.  The result is 
true in general, though, and may be proved by the same sort of reasoning, \emph{mutatis mutandis}.

As in Lemma~\ref{FGNF1}, we may assume that $G$
has the form $G(X,Y) \equiv X + Y \bmod \deg p$.   With this normalization, the 
form of the equations defining $\mathcal{A}_6 = \mathcal{A}_{4 + 2}$ is as follows.
From terms of degree $1$ in $u,v$:
\begin{align*}
u : &\sum_{i=1}^4 a_{i,0} A_{i,1}^{(j)} =0\\
v:  &\sum_{i=1}^4 A_{i,1}^{(j)} = 0
\end{align*}
From terms of degree $2$ in $u,v$:
\begin{align*}
u^2: &\sum_{i=1}^4 a_{i,0}^2 A_{i,2}^{(j)} + a_{i,0} a_{i,1} A_{i,1}^{(j)} = 0\\
uv: &\sum_{i=1}^4 2 a_{i,0} A_{i,2}^{(j)} + a_{i,1} A_{i,1}^{(j)} = 0\\
v^2: &\sum_{i=1}^4 A_{i,2}^{(j)} = 0
\end{align*}

As in item (d) in the Notation in the previous section, the equation from $u^{m - \ell} v^\ell$
is obtained from the equation for $u^m$ by applying the formal differential operator
$$\frac{1}{\ell!}\left(\frac{\partial^\ell}{\partial a_{1,0}^\ell} +  \frac{\partial^\ell}{\partial a_{2,0}^\ell} + \frac{\partial^\ell}{\partial a_{3,0}^\ell} + \frac{\partial^\ell}{\partial a_{4,0}^\ell}\right).$$
Thus the remaining equations can be obtained from the following:  from degree $3$, 
$$u^3:  \sum_{i=1}^4 a_{i,0}^3 A_{i,3} + 2 a_{i,0}^2 a_{i,1}  A_{i,2} + (a_{i,0}^2 a_{i,2} + a_{i,0} a_{i,1}^2) A_{i,1}  = 0.$$
Then from degree $4$:
\begin{align*}
u^4:  &\sum_{i=1}^4 a_{i,0}^4 A_{i,4}^{(j)} + \sum_{i=1}^4 3 a_{i,0}^3 a_{i,1} A_{i,3}^{(j)} + \sum_{i=1}^4 (2 a_{i,0}^3 a_{i, 2} + 3 a_{1, 0}^2 a_{i,1}^2) A_{i,2}^{(j)}\\
        & + (a_{i,0}^3 a_{i,3} + 3 a_{i,0}^2 a_{i,1}a_{i,2} + a_{i,0} a_{i,1}^3) A_{i,1}^{(j)} = 0.
 \end{align*}
Finally, from degrees $\ge 5$:
\begin{align*}
u^5: 0 &= \sum_{i=1}^4 a_{i,0}^5 A_{i,5}^{(j)} + \sum_{i=1}^4 4 a_{i,0}^4 a_{i,1} A_{i,4}^{(j)} + \sum_{i=1}^4 (3 a_{i,0}^4 a_{i,2} + 6 a_{i,0}^2 a_{i,1}^2) A_{i,3}^{(j)} \\
           &\qquad + \cdots  + \sum_{i=1}^4 (a_{i,0}^4 a_{i,4} + 4 a_{i,0}^3 a_{i,1} a_{i,3} + \cdots + a_{i,0} a_{i,1}^4) A_{i,1}^{(j)},
\end{align*}
and generally
\begin{align}
\label{HigherDegrees}
u^n: 0&= \sum_{i=1}^4 a_{i,0}^n A_{i,n}^{(j)} + \sum_{i=1}^4 (n-1) a_{i,0}^{n-1} a_{i,1} A_{i,n-1}^{(j)}\nonumber \\ 
& \quad + \sum_{i=1}^4 \left((n-2) a_{i,0}^{n-1} a_{i,2} + \binom{n-1}{2}  a_{i,0}^{n-2} a_{i,1}^2\right) A_{i,n-2}^{(j)} + \cdots  \\
& \quad + (a_{i,0}^{n-1} a_{i,n-1} + (n-1) a_{i,0}^{n-2} a_{i,1} a_{i,n-2 }+ \cdots  + a_{i,0} a_{i,1}^{n-1}) A_{i,1}^{(j)}.\nonumber
\end{align}

Much of the information we need is actually contained in these terms containing $A_{i,n}^{(j)}$, $A_{i,n-1}^{(j)}$ and $A_{i,1}^{(j)}$ from the terms of total degree $n$ in $u,v$.  

We begin with some observations.  We want to consider Abel relations in a $g$-dimensional formal group law with $g = (4 - 1)(4 - 2)/2 = 3$, 
so we take $1 \le j \le 3$ here.  However the form of the equations from the coefficients of $u,v, u^2, uv, v^2$ shows that there are 
$2$ degrees of freedom in the choice of the $A_{i,1}^{(j)}$ and $1$ further degree of freedom in the choice of the $A_{i,2}^{(j)}$.  Then from the terms of degree $m = 3$ in $u,v$ on
to degree $m = p - 1$ in the truncated Abel relation, the $A_{i,m}^{(j)}$  are uniquely determined by the earlier choices and the coefficients of the $C_i$.  This follows because for
each $j$ the 
coefficient matrix of the system for the $A_{i,m}^{(j)}$ have the following form.  For $m = 3$, for instance, the matrix is
$$\begin{pmatrix}
a_{1,0}^3 & a_{2,0}^3 & a_{3,0}^3 & a_{4,0}^3\\
3 a_{1,0}^2 & 3 a_{2,0}^2 & 3 a_{3,0}^2 & 3 a_{4,0}^2\\
3 a_{1,0} & 3 a_{2,0} & 3 a_{3,0} & 3 a_{4,0}\\
1 & 1 & 1 & 1
\end{pmatrix}$$
whose determinant is $9$ times the usual \emph{Vandermonde determinant} of the $a_{i,0}$.  For $m > 3$, the coefficient matrices of the 
$A_{i,\ell}^{(j)}$ terms have a similar Vandermonde form, but with more rows and numerical coefficients form the binomial coefficients $\binom{m}{\ell}$.  
We assume the $a_{i,0}$ are distinct, so the determinant of the $4 \times 4$ submatrix in the last four rows of these matrices is nonzero.  
That implies the $A_{i,m}^{(j)}$ for $m \ge 3$ are uniquely determined by the $A_{i,1}^{(j)}, A_{i,2}^{(j)}$ and the coefficients of the $C_i$.  
The upshot is that  there is a $3$-dimensional space of solutions for the $A_{i,1}$ and $A_{i,2}$.  
%This is what we meant earlier about why it might be more natural to assume the formal group law is 3-dimensional for the case $d = 4$.  

However, and this will be a key observation, the coefficients $A_{i,n}^{(j)}$ for $n \ge 3$ are always determined by exactly the same
equations in terms of the $A_{i,1}^{(j)}$, the $A_{i,2}^{(j)}$, and the coefficients $a_{i,n}$ of the equations of the $C_i$.  
In the following, it will be most convenient to make a particular choice of coordinates in $G$, so to speak, by picking one particular set of 
solutions of the equations for the $A_{i,1}^{(j)}$ and $A_{i,2}^{(j)}$.  
We will use the following notation:  For each $1 \le i \le 4$, let 
\begin{equation}
\label{deltai}
\delta_i = \prod_{j\ne i} (a_{i,0} - a_{j,0}).
\end{equation}
Then it is easy to see that the following choices give two linearly independent solutions for the equations from $u,v$ above.
\begin{equation}
\begin{matrix}
\label{Ai1s13}
A_{1,1}^{(1)} = 1/\delta_1 & A_{2,1}^{(1)}  = 1/\delta_2 &  A_{3,1}^{(1)}  = 1/\delta_3 & A_{4,1}^{(1)}  = 1/\delta_4.\\
A_{1,1}^{(3)}  = a_{1,0}/\delta_1 & A_{2,1}^{(3)} = a_{2,0}/\delta_2 &  A_{3,1}^{(3)} = a_{3,0}/\delta_3 & A_{4,1}^{(3)} = a_{4,0}/\delta_4\\
\end{matrix}
\end{equation}

Substituting these into the equations from $u^2, uv, v^2$ we can solve for the $A_{i,2}^{(1)}$, $A_{i,2}^{(3)}$.
On the other hand, with the choice $A_{i,1}^{(2)} = 0$, $1 \le i \le 4$, we have 
\begin{equation}
\begin{matrix}
\label{Ai2s2}
A_{1,2}^{(2)} = 1/(2\delta_1) & A_{2,2}^{(2)} = 1/(2\delta_2) & A_{3,2}^{(2)} = 1/(2\delta_3) & A_{4,2}^{(2)} = 1/(2\delta_4).
\end{matrix}
\end{equation}
These agree with the first terms in the truncated local expansions of the formal integrals of the differentials
\begin{equation}
\label{AbInts}
  \omega_1 = \frac{dx}{\partial f/ \partial y},\  \omega_2 = \frac{x\, dx}{\partial f/\partial y},\ \omega_3 = \frac{y\, dx}{\partial f/\partial y},\
\end{equation}
where we continue the expansion until computing the integral requires inversion of an exponent containing the prime $p$.
In other words, it is not difficult to check that in fact we can take  sets of coefficients for $j = 1,2,3$ to
 agree with the terms $\bmod\ x^p$ in the series expansions of the integrals of the Abelian differentials 
on an algebraic quartic curve in characteristic zero.  
%We will now indicate two distinct (but related) ways to establish the result of this Theorem.  As mentioned above, we only show the 
%details for $d = 4$, but the ideas extend to all $d \ge 4$ in a relatively straightforward way.  
%For the first approach, we begin by noting that the $a_{i,0}$ and $a_{i,1}$ from the $C_i$ are completely arbitrary 
%(although we require as always that the $a_{i,0}$ are distinct.  Since the $y = a_{i,0} + a_{i,1} x$ are always the 
%local expansions of algebraic quartics (e.g. the reducible quartic made up of the lines with those $4$
%equations), it follows that any $x_i \bmod\ x^2$ and the  $t_i \bmod\ x^3$ can appear in truncated Abel relations
%$\bmod \deg 3$ and furthermore the possible $A_{i,1}^{(j)}$ and $A_{i,2}^{(j)}$ always form $3$-dimensional vector
%space determined by the $a_{i,0}$ and $a_{i,1}$.  
%Furthermore, it is also true that if we consider $C_i \bmod\ x^3$ from an algebraic quartic curve, for which the 
%first formal Reiss relation
%$$a_{1,2} + a_{2,2} + a_{3,2} + a_{4,2} = 0$$
%is satisfied, there is a truncated Abel relation using $A_{i,1}^{(j)}$ and $A_{i,2}^{(j)}$ as above where the $A_{i,3}^{(j)}$
%come from the expansions $\mod \deg 4$ of the integrals in \eqref{AbIints}.  
More precisely,  letting $f(x,y) = \prod_{i=1}^4 (y - f_i(x))$, we can take
%\begin{align*}
$$t_i^{(1)}(x) \equiv  \int \left. \frac{dx}{\partial f/\partial y}\right|_{C_i}  \bmod x^p, \quad
t_i^{(2)}(x) \equiv \int \left. \frac{x\, dx}{\partial f/\partial y}\right|_{C_i} \bmod x^p, $$
and
$$t_i^{(3)}(x) \equiv \int \left. \frac{y\, dx}{\partial f/\partial y}\right|_{C_i} = \int \left. \frac{f_i(x) dx}{\partial f/\partial y}\right|_{C_i} \bmod x^p.$$
%\end{align*}
These make sense since we assume $p \ge 7$ so we can still
invert the necessary exponents to get truncated expansions for the integrals.  

%We consider the following question:  How is it possible to deform the $C_i$ to branches $\overline{C}_i$ defined by
%$$y = a_{i,0} + a_{i,0} x + (a_{i,2} + \alpha_{i,2}) x^2$$
%with corresponding $\overline{x)}(u,v)$, 
%and the $t_i^{(j)}(x)$ to mappings 
%$$\overline{t}_i^{(j)}(x) = t_i^{(j)}(x) + \mathcal{A}_{i,3}^{(j)} x^3$$
%n such a way that there is still a truncated Abel relation $\bmod \deg 4$? 

%Note that
%$$\prod_G \overline{t}_i{(j)}(\overline{x}_i(u,v)) \equiv \prod_G \overline{t}_i(x_i(u,v)) \equiv 0 \bmod \deg 3$$
%for all $j = 1,2,3$.  
%But the coefficients of $x, x^2$ in the $\overline{t_i}^{(j}(x)$ are all zero, so as a truncated Abel
%relation on the $x_i(u,v)$, this could only be extended by making all the coefficients of the higher-degree terms
%equal to $0$.
%
 %We consider the difference
%$$\prod_G \overline{t}_i (\overline{x}_i)(u,v)) - \prod_G  t_i(x_i(u,v)) = 0.$$
%Nonzero terms only start with terms of total degree 3 in $u,v$ and taking the terms from the component with $j = 1$, the form of the resulting 
%equations is
%\begin{align*}
%&\sum_{i=1}^4 a_{i,0}^3 \mathcal{A}_{i,3}^{(1)}  + \frac{1}{3} \sum_{i=1}^4 \alpha_{i,2} = 0\\
%&\sum_{i=1}^4 3 a_{i,0}^2 \mathcal{A}_{i,3}^{(1)} = 0\\
%&\sum_{i=1}^4 3 a_{i,0} \mathcal{A}_{i,3}^{(1)} = 0\\
%&\sum_{i=1}^4 \mathcal{A}_{i,3}^{(1)} = 0
%\end{align*}

% For the first approach, we consider the tangent space of $\mathcal{X}_6$ at each point and make use of a special feature of the equations related to 
%the weights considered before.  Recall that 
%by (c) of the Notation in the previous section, if ${\rm wt}(a_{i,n}) = n$ and ${\rm wt}(A_{i,n}^{(j)}) = n - 1$, then 
%these equation from the coefficient of $u^{m-\ell} v^\ell$ is isobaric of weight $m - \ell$.  If we eliminate the $A_{i,n}^{(j)}$, the resulting
%equations for 
%$$\mathcal{X}_6 \subset (\A^{20})^* = {\rm Spec}\ \ZZ[a_{i,0}, \ldots, a_{i,4};  (a_{i,0} - a_{j,0})^{-1}, 1 \le i < j \le 4]$$ 
%will also be isobaric.  It follows that if a collection of $C_i$ defined by 
%$$f_i(x) = a_{i,0} + a_{i,1} x + a_{i,2} x^2 + a_{i,3} x^3 + a_{i,4} x^4$$
%correspond to a point of $\mathcal{X}_6$, then the $C_i(t)$ defined by 
%$$f_i(t)(x) = a_{i,0} + t a_{i,1} x + t^2 a_{i,2} x^2 + t^3 a_{i,3} x^3 + t^4 a_{i,4} x^4$$
%also correspond to points of $\mathcal{X}_6$ for all $t \in k$.  Hence every 
%point of $\mathcal{X}_6$ is joined by the point corresponding to 
%$f_i(0) = a_{i,0}$ by an irreducible affine rational normal curve contained in $\mathcal{X}_6$.  
%
The $(x_i) \in \mathcal{X}_6$ from algebraic quartic curves come from a smooth, $14$-dimensional 
subvariety $\mathcal{C}_6 \subset \mathcal{X}_6$.  (Note that there is also a $\binom{4 + 2}{2} = 15$-dimensional vector space
of homogeneous polynomials of degree $4$ in $x,y$, hence a $14$-dimensional projective space of all plane quartics.)

The following fact is obvious, but it is important for our approach.
\vskip 10pt
\noindent
{\it Claim.}
The $(x_i)$ in $\mathcal{C}_6$ are the local expansions 
for which
\begin{equation}
\label{prodexp}
\prod_{i=1}^4 (y - f_i(x) \bmod x^5)
\end{equation}
has total degree $4$ in $x,y$.

\vskip 10pt
This leads to a set of equations on the coefficients in the $f_i(x)$ of the following form.  First from the coefficients of $y^3 x^2$, $y^3 x^3$, and $y^2 x^3$ in
the expansion of \eqref{prodexp},
\begin{align}
\label{Reiss4a}
0 &= \sum_{i=1}^4 a_{i,2} \nonumber \\
0 &= \sum_{i=1}^4 a_{i,3}  \\
0 &=\sum_{i=1}^4 \left(\left(\sum_{j \ne i} a_{j,0}\right) a_{i,3} + \left(\sum_{j\ne i} a_{j,1}\right) a_{i,2}\right) = 0. \nonumber
\end{align}
 In addition, from the coefficients of $y^3 x^4$, $y^2 x^4$, and $y x^4$ in \eqref{prodexp}, the following additional relations must hold on the coefficients of $x^4$:
\begin{align}
\label{Reiss4b}
0 &= \sum_{i=1}^4 a_{i,4} \nonumber \\
0 &= \sum_{i=1}^4 \left( \left(\sum_{j \ne i} a_{j,0}\right) a_{i,4} + \left(\sum_{j\ne i} a_{j,1}\right) a_{i,3} \right) + \sum_{1 \le i < j \le 4} a_{i,2} a_{j,2}  \\
0 &=\sum_{i=1}^4 \left(\left(\sum_{\substack{1 \le j < k \le 4\\  j,k \ne i}} a_{j,0} a_{k,0}\right) a_{i,4} + \sum_{i=1}^4 \left(\sum_{\substack{1 \le j,  k \le 4,\\  j,k \ne i}} a_{j,0} a_{k,1}\right) a_{i,3} \right) \nonumber\\ 
&\qquad + \sum_{i=1}^4 \left(\sum_{\substack{1 \le j < k \le 4,\\ j,k \ne i}} a_{j,1} a_{k,1} \right) a_{i,2} + \sum_{\substack{1 \le i < j \le 4,\\ \{k,\ell\} \cup \{i,j\} = \{1,2,3,4\} }} (a_{k,0} + a_{\ell,0}) a_{i,2} a_{j,2} \nonumber
\end{align}
Under the assumption that the $a_{i,0}$ are distinct, it is easy to check
that these equations define a smooth 14-dimensional subvariety 
$$\mathcal{C}_6 \subset (\A^{20})^* = {\rm Spec}\ \ZZ[a_{i,0}, \ldots, a_{i,4},  (a_{i,0} - a_{j,0})^{-1}, 1 \le i < j \le 4].$$  
We claim that in fact $\mathcal{X}_6$ equals this $\mathcal{C}_6$. 

A direct computation (we used the Maple computer algebra system for this) shows that substituting the corresponding values for the $A_{i,n}^{(1)}$, 
$n = 1, 2,3$ from 
$$t_i^{(3)}(x) =  \int \left. \frac{y\, dx}{\partial f/\partial y}\right|_{C_i} = \int \left. \frac{f_i(x) dx}{\partial f/\partial y}\right|_{C_i} \bmod x^6,$$
and simplifying, the equation from $u^3$ for $j = 3$ implies
$$0 = a_{1,2} + a_{2,2} + a_{3,2} + a_{4,2},$$
the first equation from \eqref{Reiss4a}.  (The same relation can be obtained in other ways as well from the equations using the $t_i^{(j)}$ for $j = 1,2$.)  

Similarly, the equation from $u^3 v$ for $j=3$ implies
$$0 = a_{1,3} + a_{2,3} + a_{3,3} + a_{4,3},$$
which is the second equation from \eqref{Reiss4a}.
The equation from $u^4$ for $j = 3$ implies
$$0 = \sum_{i=1}^4 (2 a_{i,0} + \sum_{j \ne i} a_{j,0}) a_{i,3} + \sum_{i=1}^4 (2 a_{i,1} + \sum_{j \ne i} a_{j,1}) a_{i,2}$$
Using the previous two equations, this clearly reduces to the form of the third equation in \eqref{Reiss4a}.  
The equation from $u^3 v^2$ for $j = 3$ implies
$$0 = a_{1,4} + a_{2,4} + a_{3,4} + a_{4,4},$$
which is the first equation from \eqref{Reiss4b}.   The equations from $u^5$ and $u^4 v$ for $j = 1$ imply other
relations on the $a_{i,n}$ with $1 \le n \le 4$.  It can be checked by a Gr\"obner basis computation that the ideal generated by the coefficients of $u^{m - \ell} v^\ell$
for all $m = 1,2,3,4,5$ is equal to the ideal generated by the six right-hand sides in \eqref{Reiss4a} and \eqref{Reiss4b}.  Hence every element of 
$\mathcal{X}_6$ comes from $C_i$ where the coefficients satisfy \eqref{Reiss4a} and \eqref{Reiss4b}, and hence agree with the local expansions
of an algebraic quartic $\bmod\ x^5$.  
\end{proof}

\end{document}
